\medskip

\begin{minipage}{0.6\linewidth}
On considère un cube ABCDEFGH
et on appelle K le milieu du segment [BC].

On se place dans le repère 
$\left(\text{A}~;~\vect{\text{AB}},~\vect{\text{AD}},~\vect{\text{AE}}\right)$ 
et on considère  le tétraèdre EFGK.

On rappelle que le volume d'un tétraèdre est donné par: 

\[V = \dfrac13 \times \mathcal{B} \times h\]

où $\mathcal{B}$ désigne l'aire d'une base et $h$ la hauteur relative à cette base.
\end{minipage} \hfill
\begin{minipage}{0.35\linewidth}
\psset{unit=1cm}
\begin{pspicture}(5.5,5.8)
\psframe(0.2,0.2)(3.7,3.7)
\psline(3.7,0.2)(5,1.9)(5,5.4)(3.7,3.7)
\psline(5,5.4)(1.5,5.4)(0.2,3.7)
\psline[linestyle=dashed](0.2,0.2)(1.5,1.9)(5,1.9)
\psline[linestyle=dashed](1.5,1.9)(1.5,5.4)
\pspolygon[linestyle=dotted,linewidth=1.25pt](0.2,3.7)(5,5.4)(4.35,1.05)
\uput[dl](0.2,0.2){\small A}\uput[dr](3.7,0.2){\small B}\uput[r](5,1.9){\small C}
\uput[dr](1.5,1.9){\small D}\uput[l](0.2,3.7){\small E}\uput[r](3.7,3.7){\small F}
\uput[ur](5,5.4){\small G}\uput[ul](1.5,5.4){\small H}\uput[dr](4.35,1.05){\small K}
\psdots(2.6,4.55)(2.275,2.375)(4.675,3.225)
\uput[u](2.6,4.55){\blue P} \uput[ur](2.275,2.375){\blue M} \uput[l](4.675,3.225){\blue N} 
\end{pspicture}
\end{minipage}

\medskip

\begin{enumerate}
\item Préciser les coordonnées des points E, F{}, G et K.

E(0~;~0~;~1), \quad F(1~;~0~;~1), \quad G(1~;~1~;~1), \quad K$\left(1~;~\frac12~;~0\right)$.
\item Montrer que le vecteur $\vect{n}\begin{pmatrix}\phantom{-}2\\-2\\\phantom{-}1\end{pmatrix}$ est orthogonal au plan (EGK).

D'après la question précédente : $\vect{\text{EG}}\begin{pmatrix}1\\1\\0 \end{pmatrix}$ et $\vect{\text{EK}}\begin{pmatrix}1\\\dfrac12\\- 1 \end{pmatrix}$.

D'où $\vect{n} \cdot \vect{\text{EG}} = 2 - 2 + 0 = 0$ : ces deux vecteurs sont orthogonaux ;

$\vect{n} \cdot \vect{\text{EK}} = 2 - 1 - 1 = 0$ : ces deux vecteurs sont orthogonaux.

Le vecteur $\vect{n}$ orthogonal à deux vecteurs non colinéaires du plan (EGK) est orthogonal à ce plan.
\item Démontrer que le plan (EGK) admet pour équation cartésienne : $2x - 2y + z - 1 = 0.$

Le résultat précédent montre qu'une équation du plan (EGK) est $2x - 2y + z = d$, avec $d \in \R$.

En particulier E$(0~;~0~;~1) \in (\text{EGK}) \iff 0 - 0 + 1 = d \iff d = 1$.

On a donc $M(x~;~y~;~z) \in (\text{EGK}) \iff 2x - 2y + z = 1$.
\item Déterminer une représentation paramétrique de la droite $(d)$ orthogonale au plan (ECK)
passant par F{}.

La droite $(d)$ étant orthogonale au plan (ECK), on peut prendre comme vecteur directeur de cette droite le vecteur $\vect{n}$. Elle contient F, donc :

$M(x~;~y~;~z) \in (d) \iff \vect{\text{F}M} = t\vect{n} \iff \left\{\begin{array}{l c l}
x - 1&=&t \times 2\\
y - 0&=&t \times (- 2)\\
z - 1&=&t \times 1
\end{array}\right., t \in \R \iff 
\left\{\begin{array}{l c l}
x &=& 1 + 2t\\
y &=&\phantom{1}- 2t\\
z &=&1 + t
\end{array}\right., t \in \R$.
\item Montrer que le projeté orthogonal L de F sur le plan (EGK) a pour coordonnées $\left(\frac59~;~\frac49~;~\frac79\right)$.

L étant le projeté orthogonal de F sur le plan (EGK), L est un point de la droite $(d)$ et du plan (EGK) : ses coordonnées vérifient donc les équations de $(d)$ et du plan (EGK), donc le système :

$\left\{\begin{array}{l c l}
x &=& 1 + 2t\\
y &=&\phantom{1}- 2t\\
z &=&1 + t\\
2x - 2y + z&=&1
\end{array}\right., t \in \R$

En remplaçant $x$, \: $y$,\: et $z$ par leurs expressions en fonction de $t$ dans la dernière équation on obtient :

$2(1 + 2t) - 2\times (- 2t) + 1 + t = 1\iff 2 + 4t + 4t +t + 1 = 1\iff 9t = - 2\iff t = - \dfrac29$.
On a donc : $x = 1 + 2\times \left(- \dfrac29\right) = 1 - \dfrac{4}{9} = \dfrac{5}{9}$ ;

$y = - 2 \times \left(- \dfrac29\right)  = \dfrac49$ ;

$z = 1 - \dfrac29 = \dfrac79$.

Conclusion : L$\left(\dfrac59~;~\dfrac49~;~\dfrac79\right)$.
\item Justifier que la longueur LF est égale à $\dfrac23$.

Avec L$\left(\frac59~;~\frac49~;~\frac79\right)$ et F(1~;~0~;~1), on a $\vect{\text{LF}}\begin{pmatrix}\frac49\\-  \frac49\\\frac29 \end{pmatrix}$.

Donc LF$^2 = \left(\frac49\right)^2 + \left(\frac49\right)^2 + \left(\frac29\right)^2  = \frac{16}{81} + \frac{16}{81} + \frac{4}{81} = \frac{36}{81}$, d'où LF $= \frac69 = \frac23$.
\item Calculer l'aire du triangle EFG. En déduire que le volume du tétraèdre EFGK est égal à $\dfrac16$.

Le triangle (EFG) est rectangle en F, donc :

$\bullet~~$On a $\mathcal{A}(\text{EFG}) = \text{EF} \times \text{FG} \times \dfrac12 = 1 \times 1 \times \frac12 = \frac12$.

$\bullet~~$$\mathcal{V}(\text{EFGK}) = \dfrac13 \times \mathcal{A}(\text{EFG}) \times \text{BF} = \frac13 \times \frac12 \times 1 = \frac16$.
\item Déduire des questions précédentes l'aire du triangle EGK.

On a aussi $\mathcal{V}(\text{EFGK}) = \dfrac13 \times \mathcal{A}(\text{EGK}) \times \text{LF}$, soit $\dfrac16 = \dfrac13 \times \mathcal{A}(\text{EGK}) \times \dfrac23$ ou encore

$\dfrac16 = \dfrac29 \times \mathcal{A}(\text{EGK})$, d'où $\mathcal{A}(\text{EGK}) = \dfrac{\frac16}{\frac29} = \dfrac16 \times \dfrac92 = \dfrac{9}{12} = \dfrac34$.
\item On considère les points P milieu du segment [EG], M milieu du segment [EK] et N milieu du segment[GK]. Déterminer le volume du tétraèdre FPMN.

D'après le théorème de la droite des milieux les côtés du triangle (PMN) ont une longueur moitié de celles du triangle (EGK). L'aire du triangle (PMN) est donc égale à $\dfrac12 \times \dfrac12 = \dfrac14$ fois celle du triangle (EGK), donc $\mathcal{A}(\text{PMN}) = \dfrac14 \times \dfrac34 = \dfrac{3}{16}$.

Comme les trois points P, M et N appartiennent au plan (EGK) le triangle (PMN) appartient au plan (EGK) : la hauteur du tétraèdre FPMN est donc la même que celle du tétraèdre (EFGK) soit LF.

On a donc $\mathcal{V}(\text{FPMN}) = \dfrac13 \times \mathcal{A}(\text{PMN}) \times \text{LF} = \dfrac13 \times \dfrac{3}{16} \times \dfrac{2}{3} = \dfrac{1}{24}$.
\end{enumerate}

\bigskip
